% ===========================================================================
%  Method section for Reasoning Speculation (NeurIPS 2026)
% ===========================================================================

We now present \method{} (\fullmethod{}), a training-free, three-phase
framework that restructures test-time reasoning compute by using
\emph{parallel short explorations} to estimate problem difficulty via
cross-path consensus, then adaptively routing each question to an
appropriately sized resolution strategy.

% ---------------------------------------------------------------------------
\subsection{Problem Setup and Notation}
\label{sec:method:setup}
% ---------------------------------------------------------------------------

Let $\mathcal{M}$ denote a reasoning-capable language model (e.g.\ one with
a chain-of-thought or ``thinking'' mode) and let $q$ be an input question
drawn from a benchmark $\mathcal{Q}$.
A \emph{reasoning trace} $y = \mathcal{M}(q, b)$ is the text generated by
$\mathcal{M}$ when prompted with $q$ under a token budget~$b$.
We write $a(y)$ for the answer extracted from $y$ by a deterministic parser
(regex-based; see Appendix for details) and $|y|$ for the number of tokens
consumed.

We consider three budget levels
$B_{\mathrm{probe}} < B_{\mathrm{med}} < B_{\mathrm{hard}}$
(e.g.\ 128, 256, 512 tokens) and a parallelism factor $K \geq 2$.
The goal is an \emph{adaptive policy} $\pi$ that, given $q$ and
$\mathcal{M}$, selects an execution plan maximizing the cost-aware utility
\begin{equation}
  \label{eq:rs-utility}
  U(\pi) \;=\; \mathbb{E}_{q}\!\Big[\,
    \mathbf{1}\big[\,a(\pi(q)) = a^{\star}(q)\,\big]
    \;-\; \lambda \cdot \frac{|\pi(q)|}{B_{\mathrm{hard}}}
  \Big],
\end{equation}
where $a^{\star}(q)$ is the gold answer and $\lambda \geq 0$ trades off
accuracy against token cost.
A \emph{fixed-budget} baseline simply sets $\pi(q) = \mathcal{M}(q, B)$
for a constant~$B$; \emph{self-consistency}~(SC@$K$) samples $K$ full-budget
traces and majority-votes.
Both are compute-agnostic---they allocate the same resources regardless of
problem difficulty.

% ---------------------------------------------------------------------------
\subsection{Phase~1: Parallel Exploration}
\label{sec:method:explore}
% ---------------------------------------------------------------------------

Given question $q$, \method{} generates $K$ short \emph{exploration paths}
$\{y_1, \ldots, y_K\}$, each constrained to $B_{\mathrm{probe}}$ tokens:
\begin{equation}
  \label{eq:explore}
  y_k = \mathcal{M}(q,\, B_{\mathrm{probe}},\, \tau_k),
  \quad
  \tau_k =
  \begin{cases}
    0   & k = 1 \;\text{(greedy)}, \\
    \tau & k \geq 2 \;\text{(stochastic, } \tau > 0\text{)}.
  \end{cases}
\end{equation}
The first path uses greedy decoding ($\tau_1 = 0$) to obtain a high-quality
deterministic baseline; the remaining $K{-}1$ paths use temperature
$\tau > 0$ (default 0.7) to encourage diverse reasoning strategies.
This design mirrors the ``draft'' phase of speculative
decoding~\citep{leviathan2023fast,chen2023accelerating}: it produces
lightweight probes whose primary purpose is to \emph{estimate difficulty},
not to solve the problem outright.

\paragraph{Answer extraction and projection.}
From each path $y_k$, we extract a candidate answer $a_k = a(y_k)$ via
deterministic parsing.
Because $B_{\mathrm{probe}}$ may be insufficient for the model to reach a
``Final answer: \ldots'' marker, some paths yield $a_k = \varnothing$.
For these, we apply a lightweight \emph{projection} step: a short prompt
concatenates $q$ and the partial trace $y_k$ and asks the model to output a
single answer line, consuming at most 16 additional tokens.
Projection rescues partial reasoning that is substantively correct but
merely truncated before an answer marker, and adds negligible cost.

% ---------------------------------------------------------------------------
\subsection{Phase~2: Consensus-Based Difficulty Estimation}
\label{sec:method:decide}
% ---------------------------------------------------------------------------

The key insight underlying \method{} is that \emph{multi-path agreement is a
far more reliable difficulty signal than any feature extractable from a
single reasoning trace}.
Prior work~\citep{wang2023selfconsistency} observed that self-consistency
correlates with correctness; we operationalize this as a three-component
\emph{consensus signal} that drives an explicit routing decision.

\paragraph{Consensus signal.}
Let $\mathcal{A} = \{a_k : a_k \neq \varnothing\}$ be the set of valid
extracted answers and let $a^{\mathrm{maj}}$ be the answer with the highest
frequency in $\mathcal{A}$ (ties broken arbitrarily).
We define:
\begin{align}
  \label{eq:agreement}
  \alpha &= \frac{|\{a_k \in \mathcal{A} : a_k = a^{\mathrm{maj}}\}|}
                  {|\mathcal{A}|},
  &\quad&\text{(agreement ratio)}\\[4pt]
  \label{eq:validity}
  \nu &= \frac{|\mathcal{A}|}{K},
  &\quad&\text{(prediction validity)}\\[4pt]
  \label{eq:marker}
  \mu &= \frac{1}{K}\sum_{k=1}^{K}
         \mathbf{1}[\text{path } y_k \text{ contains a ``Final answer'' marker}],
  &\quad&\text{(final-marker rate)}
\end{align}
and an overall confidence score
\begin{equation}
  \label{eq:confidence}
  \gamma = 0.5\,\alpha + 0.3\,\nu + 0.2\,\mu \;\;\in [0, 1].
\end{equation}
Agreement $\alpha$ receives the largest weight because it directly measures
whether independent reasoning chains converge to the same answer.
Validity $\nu$ captures whether the probe budget was sufficient for the model
to produce any answer at all, and marker rate $\mu$ indicates whether traces
reached a natural conclusion.

\paragraph{Routing rule.}
The consensus signal determines the difficulty route:
\begin{equation}
  \label{eq:routing}
  \mathrm{route}(q) =
  \begin{cases}
    \textsc{Easy}
      & \text{if } \alpha \geq \theta_{\mathrm{E}} \;\text{and}\; \gamma \geq 0.6, \\[3pt]
    \textsc{Medium}
      & \text{if } \alpha \geq \theta_{\mathrm{M}} \;\text{or}\; \gamma \geq 0.4, \\[3pt]
    \textsc{Hard}
      & \text{otherwise},
  \end{cases}
\end{equation}
where $\theta_{\mathrm{E}}$ and $\theta_{\mathrm{M}}$ are agreement
thresholds (default 0.75 and 0.5, respectively).
The routing is deliberately conservative: a question is labeled \textsc{Easy}
only when a super-majority of paths agree \emph{and} overall confidence is
high; the \textsc{Hard} route is reserved for genuine disagreement.

% ---------------------------------------------------------------------------
\subsection{Phase~3: Adaptive Resolution}
\label{sec:method:resolve}
% ---------------------------------------------------------------------------

Each route prescribes a resolution strategy with a distinct token budget:

\paragraph{\textsc{Easy} route\,---\,majority vote.}
When paths exhibit strong consensus, the majority answer $a^{\mathrm{maj}}$
is returned immediately.
No additional generation is performed.

\paragraph{\textsc{Medium} route\,---\,extend the best path.}
Among exploration paths whose answer matches $a^{\mathrm{maj}}$, we select
the best path $y^{\star}$ (defaulting to the greedy path $y_1$) and
\emph{extend} its reasoning by continuing generation from $y^{\star}$ up to
$B_{\mathrm{med}}$ total tokens.
If the extension reveals a revised answer, we adopt it; otherwise, the
original answer from $y^{\star}$ is retained.
Concretely, the remaining budget is $B_{\mathrm{med}} - |y^{\star}|$
tokens, and extension uses greedy decoding ($\tau = 0$) to maximize
reliability.

\paragraph{\textsc{Hard} route\,---\,full deliberation.}
For problems with low consensus, we construct a \emph{deliberation prompt}
that includes the original question $q$ and summaries of all $K$ exploration
paths (truncated to avoid exceeding the context window).
The model then generates a fresh solution of up to $B_{\mathrm{hard}}$
tokens, informed by the diverse partial solutions as context:
\begin{equation}
  y_{\mathrm{delib}} = \mathcal{M}\!\big(\,
    [q;\; \mathrm{summarize}(y_1, \ldots, y_K)],\;
    B_{\mathrm{hard}},\; \tau = 0
  \big).
\end{equation}
This deliberation is analogous to an ``ensemble distillation'' within a
single model: the diverse reasoning strategies surfaced during exploration
serve as scaffolding for a more thorough solution.

% ---------------------------------------------------------------------------
\subsection{Full Algorithm}
\label{sec:method:algorithm}
% ---------------------------------------------------------------------------

Algorithm~\ref{alg:reasonspec} summarizes the complete \method{} pipeline.

\begin{algorithm}[t]
\caption{\method{}: Reasoning Speculation}
\label{alg:reasonspec}
\begin{algorithmic}[1]
\REQUIRE Question $q$, model $\mathcal{M}$, parallelism $K$,
  budgets $(B_{\mathrm{probe}}, B_{\mathrm{med}}, B_{\mathrm{hard}})$,
  thresholds $(\theta_{\mathrm{E}}, \theta_{\mathrm{M}})$,
  temperature $\tau$
\ENSURE Answer $\hat{a}$, tokens consumed $T$

\medskip
\STATE \textbf{// Phase 1: Parallel Exploration}
\FOR{$k = 1$ \textbf{to} $K$}
  \STATE $\tau_k \leftarrow 0$ if $k{=}1$ else $\tau$
    \hfill\COMMENT{Greedy first, diverse rest}
  \STATE $y_k \leftarrow \mathcal{M}(q,\; B_{\mathrm{probe}},\; \tau_k)$
  \STATE $a_k \leftarrow \textsc{Parse}(y_k)$
  \IF{$a_k = \varnothing$}
    \STATE $a_k \leftarrow \textsc{Project}(q, y_k)$
      \hfill\COMMENT{Lightweight projection, $\leq 16$ tokens}
  \ENDIF
\ENDFOR

\medskip
\STATE \textbf{// Phase 2: Consensus \& Routing}
\STATE Compute $\alpha, \nu, \mu, \gamma$ via
  Eqs.~\eqref{eq:agreement}--\eqref{eq:confidence}
\STATE $\mathrm{route} \leftarrow$ apply Eq.~\eqref{eq:routing}
\STATE $T \leftarrow \sum_{k=1}^{K} |y_k|$
  \hfill\COMMENT{Exploration cost}

\medskip
\STATE \textbf{// Phase 3: Adaptive Resolution}
\IF{$\mathrm{route} = \textsc{Easy}$}
  \STATE $\hat{a} \leftarrow a^{\mathrm{maj}}$
    \hfill\COMMENT{Return majority vote; no extra tokens}
\ELSIF{$\mathrm{route} = \textsc{Medium}$}
  \STATE $y^{\star} \leftarrow$ best path matching $a^{\mathrm{maj}}$
  \STATE $y_{\mathrm{ext}} \leftarrow \mathcal{M}(
    [q;\, y^{\star}],\; B_{\mathrm{med}} - |y^{\star}|,\; 0)$
  \STATE $\hat{a} \leftarrow \textsc{Parse}(y^{\star} \!\oplus\! y_{\mathrm{ext}})$;
    \;$T \leftarrow T + |y_{\mathrm{ext}}|$
\ELSE \hfill\COMMENT{$\textsc{Hard}$}
  \STATE $p \leftarrow [q;\; \textsc{Summarize}(y_1, \ldots, y_K)]$
  \STATE $y_{\mathrm{delib}} \leftarrow \mathcal{M}(p,\;
    B_{\mathrm{hard}},\; 0)$
  \STATE $\hat{a} \leftarrow \textsc{Parse}(y_{\mathrm{delib}})$;
    \;$T \leftarrow T + |y_{\mathrm{delib}}|$
\ENDIF
\RETURN $\hat{a},\; T$
\end{algorithmic}
\end{algorithm}

% ---------------------------------------------------------------------------
\subsection{Complexity Analysis}
\label{sec:method:complexity}
% ---------------------------------------------------------------------------

\paragraph{Token cost by route.}
Let $C_{\mathrm{explore}} = K \cdot B_{\mathrm{probe}}$ denote the
worst-case exploration cost (actual cost may be lower when paths terminate
early).
The total token budget consumed per question is:
\begin{equation}
  \label{eq:cost}
  T(q) =
  \begin{cases}
    C_{\mathrm{explore}}
      & \textsc{Easy}, \\[2pt]
    C_{\mathrm{explore}} + (B_{\mathrm{med}} - |y^{\star}|)
      & \textsc{Medium}, \\[2pt]
    C_{\mathrm{explore}} + B_{\mathrm{hard}}
      & \textsc{Hard}.
  \end{cases}
\end{equation}
With default parameters ($K{=}4$, $B_{\mathrm{probe}}{=}128$,
$B_{\mathrm{med}}{=}256$, $B_{\mathrm{hard}}{=}512$), the three routes
cost at most 512, 640, and 1024 tokens respectively.
If a fraction $p_E$ of questions are routed as \textsc{Easy}, $p_M$ as
\textsc{Medium}, and $p_H$ as \textsc{Hard}, the expected cost satisfies
\begin{equation}
  \label{eq:expected-cost}
  \mathbb{E}[T] = C_{\mathrm{explore}}
    + p_M \cdot \bar{E}_{\mathrm{med}}
    + p_H \cdot B_{\mathrm{hard}},
\end{equation}
where $\bar{E}_{\mathrm{med}}$ is the average extension length for
\textsc{Medium} questions.
When the majority of questions are easy (as is typical for well-studied
benchmarks), the expected cost approaches $K \cdot B_{\mathrm{probe}}$,
which can be substantially below a single full-budget pass at
$B_{\mathrm{hard}}$.

\paragraph{Latency.}
The $K$ exploration paths are \emph{embarrassingly parallel}: with
sufficient hardware, Phase~1 completes in the wall-clock time of a single
$B_{\mathrm{probe}}$-token generation, not $K$ times that.
This is a practical advantage over sequential methods such as progressive
halting or iterative refinement.
The resolution phase is a single sequential generation, so the total
wall-clock latency (with $K$-way parallelism) is bounded by
$B_{\mathrm{probe}} + \max(0, B_{\mathrm{med}}, B_{\mathrm{hard}})$
tokens, comparable to a single full-budget pass for \textsc{Medium}
and \textsc{Hard} questions, and strictly less for \textsc{Easy} ones.

% ---------------------------------------------------------------------------
\subsection{Comparison with Related Approaches}
\label{sec:method:comparison}
% ---------------------------------------------------------------------------

\method{} shares components with several existing paradigms but differs in
how it combines them.

\paragraph{Self-consistency (SC).}
SC@$K$~\citep{wang2023selfconsistency} generates $K$ full-budget
traces and majority-votes.
\method{} also uses majority voting, but \emph{only as a subroutine for
easy questions}; critically, it applies voting to \emph{short} probes
($B_{\mathrm{probe}}$ tokens) rather than full-length traces, and uses the
agreement signal to \emph{decide} whether full-length generation is
warranted.
SC pays $K \times B_{\mathrm{full}}$ tokens unconditionally; \method{}
pays $K \times B_{\mathrm{probe}}$ for easy questions and invests
additional compute only where consensus is low.
When the fraction of easy questions is non-trivial, this yields large
savings: for instance, with $B_{\mathrm{probe}} / B_{\mathrm{hard}} = 1/4$
and 50\% easy questions, \method{} uses roughly $3/8$ the tokens of SC@$K$
at comparable accuracy.

\paragraph{Best-of-$N$ (BoN) and reward-model reranking.}
BoN generates $N$ complete solutions and selects the one with the highest
reward model score~\citep{cobbe2021training,lightman2024lets}.
This requires (i)~$N$ full-budget generations and (ii)~a trained reward or
verifier model.
\method{} requires neither: it generates only short probes, needs no
external verifier, and uses consensus as a zero-shot quality signal.
When consensus is insufficient (\textsc{Hard} route), \method{} invests in
\emph{a single deeper generation} informed by all probes, rather than hoping
that one of $N$ independent samples happens to be correct.

\paragraph{Monte Carlo Tree Search (MCTS) for reasoning.}
MCTS-based methods~\citep{hao2023reasoning,xie2024selfevaluation} build a
search tree over reasoning steps, using a value function to guide
expansion.
While powerful, they require step-level value estimation (often from a
separate model), multiple sequential rollouts, and careful tuning of
exploration constants.
\method{} operates at the \emph{trace level} rather than the step level,
needs no value model, and makes a single routing decision after a fixed
number of parallel probes---resulting in a simpler, more predictable
compute profile.

\paragraph{Analogy to speculative decoding.}
The name \fullmethod{} is chosen deliberately: just as speculative
decoding~\citep{leviathan2023fast,chen2023accelerating} changes the
\emph{compute structure} of token generation (small-model draft $\to$
large-model verification) without altering the output distribution,
\method{} changes the \emph{compute structure of reasoning} (parallel
short probes $\to$ consensus $\to$ adaptive-depth resolution) without
modifying the model itself.
Both methods exploit a structural asymmetry---most tokens/problems are easy
and can be handled cheaply---to reduce the expected cost while preserving
or improving quality.
